\section*{Penne all'arrabbiata}
\ihead{}\chead{Personen:4}\ohead{}
\ifoot{Vorbereitungszeit:45 min}\cfoot{}\ofoot{Kochzeit:20 min}
{\Large Zutaten}
\begin{multicols}{2}
\begin{itemize}
\item \SI{400}{g} Penne
\item \SI{100}{g} durchwachsener Speck
\item \SI{500}{g} Tomaten, frisch oder \num{1} große Dose Pelati
\item \num{1} Bund glatte Petersilie
\item \num{2} kleine rote Chilischoten
\item \SI{50}{g} Parmesan oder Pecorino, frisch gerieben
\item \num{1} Zwiebel
\item \num{2} Knoblauchzehen
\item \SI{2}{EL} Butter
\item Salz und frisch gemahlener Pfeffer
\end{itemize}
% \columnbreak
\end{multicols}
\noindent {\Large Zubereitung}\\
\\
Den Speck in feine Streifen schneiden.
Die Tomaten enthäuten, entkernen, das Fruchtfleisch klein schneiden und durch ein Sieb streichen.
Die Petersilie und die Zwiebel fein hacken.
Die Knoblauchzehen in feine Scheiben schneiden.
In einem Topf \SI{4}{l} Salzwasser aufkochen und die Penne darin \SI{5}{min} vorgaren.
In einer großen Pfanne die Butter erwärmen.
Speck und Zwiebeln hineingeben und auf kleiner Hitze unter Rühren dünsten, ohne Farbe nehmen zu lassen.
Knoblauch, passierte Tomaten und Chilischoten unterrühren und mit Salz und Pfeffer würzen.
Auf kleiner Hitze weiter köcheln lassen.
Die Penne abgießen und einige EL des heißen Kochwassers zurückbehalten.
Die Penne nur kurz abtropfen lassen und unter die Sauce mischen.
Auf kleiner Flamme weiter köcheln lassen, bis die Nudeln al dente sind.
Wenn nötig, etwas Kochwasser nachgießen.
Eventuell die Chilischoten entfernen.
Nochmals mit Salz und Pfeffer abschmecken.
Je nach Belieben den Käse unter die Nudeln mischen und in einer vorgewärmten Schüssel oder auf tiefen Tellern anrichten.
